% !TEX program = xelatex
\documentclass[11pt,a4paper]{article}

\usepackage[margin=1.1cm, top=0.8cm, bottom=0.8cm]{geometry}
\usepackage{fontspec}
\usepackage{xeCJK}
\usepackage[hidelinks,unicode]{hyperref}
\usepackage{enumitem}
\usepackage{titlesec}
\usepackage{tabularx}
\usepackage{xcolor}
\usepackage{fontawesome5}
\usepackage{setspace}

\pagestyle{empty}
\definecolor{mainblue}{RGB}{25,80,150}
\setmainfont{Helvetica Neue}[Ligatures=TeX]
\setmonofont{Menlo}
\setCJKmainfont{FandolSong-Regular.otf}
\setCJKsansfont{FandolHei-Regular.otf}
\renewcommand{\familydefault}{\sfdefault}
\setlist[itemize]{noitemsep, topsep=0pt, left=0pt, itemsep=2pt, parsep=0pt, labelsep=0.5em, leftmargin=1em}
\titleformat{\section}{\large\bfseries\color{mainblue}}{}{0em}{}[\color{mainblue}\titlerule]
\titlespacing*{\section}{0pt}{2.5ex}{1ex}

\newcommand{\projecttitle}[1]{\textbf{#1}}
\newcommand{\techstack}[1]{\textit{技术栈: #1}}
\newcommand{\projectdesc}[1]{\textbf{项目简介：}#1}

\begin{document}
\setstretch{0.92}

\begin{center}
    {\Huge\bfseries 你的姓名} \\[1em]
    {\large Java 后端 / 求职方向} \\[0.5em]
\end{center}

\section*{联系方式}
\begin{tabularx}{\textwidth}{X X X X}
    \faPhone\ 电话: 你的电话 &
    \faEnvelope\ 邮箱: your@email.com &
    \faGithub\ GitHub: yourname &
    \faGlobe\ 个人网站: example.com
\end{tabularx}

\section*{教育经历}
\begin{tabularx}{\textwidth}{@{}XXX@{}}
  你的学校 & 你的专业 & 2023.09 -- 2027.06
\end{tabularx}

\section*{实习经历}
\projecttitle{某科技公司 - 示例项目} \hfill \\
\techstack{Spring Boot, MySQL, Redis}

\projectdesc{一句话说明业务背景、你的职责和项目价值。}

\begin{itemize}
    \item \textbf{模块开发：} 描述你负责的核心模块、接口或微服务，以及你承担的职责边界。
    \item \textbf{性能优化：} 描述你如何处理慢 SQL、缓存、异步任务、并发一致性或系统稳定性问题。
    \item \textbf{结果表达：} 尽量补充真实结果、收益或可复盘的问题定位过程。
\end{itemize}

\section*{个人项目}
\projecttitle{示例项目名称} \hfill \\
\techstack{Spring Boot, PostgreSQL, Redis, RabbitMQ}

\projectdesc{一句话说明项目功能和核心亮点。}

\begin{itemize}
    \item \textbf{权限设计：} 描述登录、鉴权、RBAC 或资源级权限控制设计。
    \item \textbf{业务设计：} 描述核心业务流程、关键接口或数据模型设计。
    \item \textbf{稳定性优化：} 描述缓存、消息队列、幂等、限流、重试或监控相关实践。
\end{itemize}

\section*{技术能力}
\begin{itemize}
    \item \textbf{Java 基础：} 按你真实掌握情况填写集合、多线程、JVM、并发工具等。
    \item \textbf{Spring 生态：} 按你真实掌握情况填写 Spring Boot、Spring MVC、Spring Security 等。
    \item \textbf{数据库与中间件：} 按你真实掌握情况填写 MySQL、Redis、RabbitMQ 等。
    \item \textbf{工程化能力：} 按你真实掌握情况填写 Git、Docker、Linux 等。
\end{itemize}

\end{document}
